\topic{15}{Off-design robustness at Mach 1.7 and 1.9}
The surrogate training campaign fixes Mach number at 1.8. To test the \emph{geometry}, rather than asking the network to extrapolate in Mach number, the baseline and retained candidate were solved directly with CFD at $M=1.7$ and $M=1.9$.

At $M=1.7$, the optimized geometry reduces peak pressure by 21.01\% and pressure drag by 3.23\%. At $M=1.9$, the peak reduction is 21.17\% and drag reduction is 4.31\%. These are direct CFD stress tests of the fixed geometry. They do not establish robustness over a complete flight envelope.

\keyidea{A geometry can be tested outside the training operating condition without pretending that the neural network was trained to predict those new conditions.}

\topic{16}{NASA SEEB-ALR: an independent reference, not a successful Week 16 validation}
The NASA SEEB-ALR benchmark serves a different purpose from the Taylor--Maccoll cone and from the two-parameter learning body. The public reference includes as-built geometry, wind-tunnel pressure records, coordinate-alignment information, and a NASA-hosted LAVA CFD solution. The comparison condition is $M=1.6$ at zero incidence.

NASA reports pressure excess normalized by $p_\infty$, while our design-body plots use $\Cp$. Section 3 showed the required conversion. The experimental curves also carry coordinate shifts defined by the original processing macros; those shifts must be preserved rather than tuned to align shocks visually.

\begin{figure}[H]\centering
\includegraphics[width=.92\linewidth]{seeb_comparison.png}
\caption{Retained NASA SEEB-ALR reference comparison. The LAVA result is an independent numerical solution, not an exact answer. Experimental agreement and solver-to-solver agreement answer different questions.}
\end{figure}

\topic{17}{The failed NASA reproduction is scientifically useful}
The attempted SU2 NASA replay did not establish an accepted mesh-refined, physically admissible family. Earlier nose treatment produced unphysical stagnation-region behavior; later recovery runs still failed the declared convergence/physical gates. Those failures are retained rather than hidden.

This is an important lesson for scientific machine learning: a visually plausible pressure trace or a solver exit code is not enough. A reference CFD label must first survive geometry, numerical, and physical checks before it can be used to validate a learned model.

In the release record, NASA is explicitly marked \texttt{failed/deferred}; the accepted Week 16 core rests on the 44-case teaching-body campaign, the Taylor--Maccoll verification, the clean neural model, and the independently recomputed retained design.

\checkpoint{Would agreement between SU2 and LAVA, without comparison to experiment, constitute experimental validation? Why not?}

\topic{18}{What is still required before claiming a quieter ground boom?}
A complete low-boom prediction chain would require, at minimum:
\begin{enumerate}
\item a realistic aircraft geometry with volume and lift distributions, including wings and propulsion integration;
\item near-field CFD verified for mesh, numerical convergence, and physical consistency;
\item validation against an appropriate measured near-field dataset;
\item transfer of the near-field signature to a validated propagation solver;
\item atmospheric profiles for temperature, humidity, wind, and molecular relaxation;
\item ground-waveform prediction under several operating/atmospheric conditions;
\item calculation of an accepted psychoacoustic metric such as PLdB with its assumptions and uncertainty.
\end{enumerate}
Consequently,
\begin{equation}
 \boxed{\text{lower near-field peak}\;\not\Rightarrow\;\text{automatically lower ground PLdB}.}
\end{equation}
The Week 16 result is intentionally narrower: within the verified axisymmetric Euler teaching problem, the retained design reduces the off-body peak pressure while satisfying the pressure-drag constraint on both design meshes.

\topic{19}{Reproducibility, provenance, and final takeaways}
The clean campaign uses SU2 8.0.1 and contains 44 accepted cases. Model fitting uses only the 24 training geometries, 12 full-SVD POD modes, and a fixed 32--32 tanh MLP. The retained design is not accepted on surrogate prediction alone; it is recomputed with CFD on two meshes and at two nearby Mach numbers. Raw/compact evidence, hashes, and audit reports are retained separately from classroom-friendly arrays.

\textbf{The scientific workflow to remember is}
\begin{equation}
 \boxed{
 \text{physics}\rightarrow\text{verified CFD}\rightarrow\text{POD}\rightarrow
 \text{surrogate}\rightarrow\text{optimization}\rightarrow\text{fresh CFD verification}
 }.
\end{equation}

\textbf{Five final lessons.}
\begin{enumerate}
\item A small residual is not the same as an accurate or physically admissible CFD solution.
\item A small average neural error can hide a poor individual waveform.
\item Good interpolation does not imply safe extrapolation.
\item Optimizers exploit model error; recompute the proposed optimum with the governing equations.
\item Near-field pressure shaping and ground sonic-boom loudness are distinct stages of the physics.
\end{enumerate}

\textbf{Questions for discussion.}
\begin{enumerate}
\item Why is fixed volume necessary for a fair geometry optimization?
\item What does a POD mode represent physically, and why must it be fitted on the training set only?
\item Why can the peak error be small while waveform $L_2$ error is large?
\item Which evidence in this module is verification, which is ML validation, and which is experimental reference evidence?
\item What would you add first if the goal changed from ``reduce near-field peak'' to ``minimize ground PLdB''?
\end{enumerate}

